% ============================================================
% 3.4.3 RISK CONTROL
% ============================================================

\hypertarget{p3-4-3}{}

\subsection{Risk Control}

\subsubsection{General}

Risk control for the \textbf{BN-100 LUX} Infant Radiant Warmer is addressed
through the risk-control measures documented by the manufacturer.

The identified control measures are related to the design, protective
functions, monitoring systems, alarms, components and manufacturing
processes of the device.

The available Risk Management Report constitutes the primary documentary
basis for the risk-control measures presented in this section.

\subsubsection{Risk Control Principles}

The manufacturer implements risk-control measures through the design and
operation of the device.

The documented controls include, where applicable:

\begin{itemize}

    \item inherently safe design characteristics;

    \item protective measures;

    \item monitoring functions;

    \item alarm functions;

    \item electrical protection;

    \item thermal protection;

    \item temperature monitoring;

    \item mechanical limitations;

    \item manufacturing and process controls.

\end{itemize}

The risk-control measures are considered in relation to the identified
hazards and hazardous situations.

\subsubsection{Electrical Risk Controls}

Electrical safety is addressed through the use of documented electrical
components and protective elements.

The manufacturer identifies safety-related components including:

\begin{itemize}

    \item dual fuse input socket;

    \item power supply cord;

    \item fuse;

    \item relay;

    \item rocker switch;

    \item switching power supply;

    \item filter;

    \item internal electrical wiring;

    \item printed circuit board.

\end{itemize}

These components contribute to the control of electrical hazards associated
with the operation of the device.

\subsubsection{Thermal Risk Controls}

Thermal safety is addressed through temperature control, temperature
monitoring and over-temperature protection.

The manufacturer documentation specifies:

\begin{itemize}

    \item controlled temperature range of $32$--$38\,^\circ$C;

    \item over-temperature alarm at the documented thresholds;

    \item heating cut-off under specified over-temperature conditions;

    \item temperature deviation alarm;

    \item skin temperature monitoring.

\end{itemize}

These functions provide documented controls against hazardous temperature
conditions.

\subsubsection{Temperature Sensor Controls}

The skin temperature sensor is part of the temperature monitoring and
control system.

The manufacturer documents alarm functions for sensor failure conditions,
including:

\begin{itemize}

    \item skin temperature sensor short circuit;

    \item skin temperature sensor open circuit.

\end{itemize}

These alarm functions provide a means of detecting specified sensor faults
and contribute to the control of risks associated with incorrect or
unavailable temperature measurements.

\subsubsection{Alarm Controls}

The device incorporates audible and visual alarm functions associated with
specified abnormal conditions.

The documented alarm functions include:

\begin{itemize}

    \item over-temperature alarm;

    \item temperature deviation alarm;

    \item skin temperature sensor short-circuit alarm;

    \item skin temperature sensor open-circuit alarm;

    \item power failure alarm;

    \item manual-mode alarm.

\end{itemize}

The alarms provide information to the operator concerning specified
abnormal operating conditions.

\subsubsection{Mechanical Risk Controls}

Mechanical risks are controlled through defined mechanical characteristics
and operating limits.

The manufacturer specifies limits for components including:

\begin{itemize}

    \item mattress support;

    \item tray;

    \item crib inclination and adjustment.

\end{itemize}

The documented load limits provide defined conditions for the intended use
of the corresponding components.

\subsubsection{Patient-Contact Risk Controls}

The manufacturer identifies patient-contact components within the risk
management documentation.

These include:

\begin{itemize}

    \item cotton baby mattress;

    \item stainless-steel skin temperature sensor probe.

\end{itemize}

The corresponding biological considerations are addressed through the
manufacturer's risk management process.

No additional biological safety claim is introduced beyond the available
manufacturer documentation.

\subsubsection{Manufacturing and Process Controls}

The manufacturer's risk management documentation considers manufacturing
and assembly processes as part of the overall risk-control framework.

Controls are applied to relevant manufacturing steps and critical
components to reduce the possibility that manufacturing defects could
adversely affect the safety or performance of the finished device.

The available documentation does not provide a complete independent
validation package for every manufacturing process.

Accordingly, no additional process-validation claim is introduced in this
section.

\subsubsection{Information for Safety}

The user manual provides information necessary for the appropriate operation
of the device, including operating conditions, temperature-control modes,
alarm conditions and applicable limitations.

The provision of operating instructions and warnings forms part of the
documented risk-control framework.

The operator is expected to use the device in accordance with the
manufacturer's instructions.

\subsubsection{Risk-Control Verification}

The effectiveness of risk-control measures is addressed through the
manufacturer's technical documentation and risk management process.

Where objective verification requires a test report, inspection record or
other evidence, the corresponding manufacturer documentation should be
included in the technical file.

No independent test result is introduced in this section where such
evidence is not available in the submitted documentation.

\subsubsection{Residual Risk}

Following implementation of the documented risk-control measures, the
remaining risks are considered within the manufacturer's overall risk
management process.

The assessment of the remaining risks is addressed in Section 3.4.4,
``Overall Residual Risk Evaluation''.

\subsubsection{Documentation Limitation}

This section summarizes the risk-control measures documented by the
manufacturer.

No new risk-control measure, acceptance criterion or performance claim is
introduced independently of the available manufacturer documentation.

\subsubsection{Documentary Basis}

This section is based primarily on:

\begin{itemize}

    \item Risk Management Report (BN-100/BN-100A), Document No.
    YF-7.3-24-02 (w), Version A/0;

    \item BN-100 / BN-100 LUX User Manual.

\end{itemize}
